\section*{Abstract}
Current intersection management strategies rely heavily on signalized controllers despite causing delay, increased energy consumption, and road accidents. However, the emergence of autonomous vehicles (AVs) offers an opportunity for more efficient intersection management strategies, such as a communication-free distributed barrier control algorithm. This thesis evaluates how such a controller performs against a Fully Actuated System (FAS) controller at a single intersection, using the CARLA physics based simulator. The barrier and FAS controller are simulated on a closed-loop figure-eight track and compared across four traffic densities (\(N= 51,61,71,81\) vehicles) over 30 simulations under varying initial conditions. Performance is evaluated in terms of throughput, delay time, energy consumption, and safety violations. Results highlight that the barrier controller outperforms the FAS controller across all traffic densities in throughput, delay time, and energy consumption. The out performance is attributed to the forming of a steady-state platoon which allows for continuous intersection crossing without stop-and-go behavior. However, at higher traffic densities, \(N=71\) and \(N=81\), safety violations were recorded using the barrier controller, whereas the FAS controller recorded near zero violations. The violations were only recorded during the phase before the steady-state platoon was formed, which indicates that the barrier controller itself is capable of safe performance at higher traffic densities. In conclusion, the barrier controller is found to be a promising solution for intersection management for autonomous vehicles, with room  for improvement in terms of barrier parameters to ensure completely safe operation.